\section{Parameter dependence and limiting cases}

\Cref{eq:Hpsm,eq:ellipse-center,eq:semiaxes} separate parameter effects.

\paragraph{$R_s\to0$.}
The cross term vanishes and the axes align with $d$ and $q$:
\begin{equation}
 \frac{(\id+\psif/\Ld)^2}{(\Umax/\abs{\omegae}\Ld)^2}
 +\frac{\iq^2}{(\Umax/\abs{\omegae}\Lq)^2}=1.
\end{equation}
Resistance is therefore the source of the small rotation in the linear model.

\paragraph{$L_d=L_q=L_s$.}
The shape matrix becomes
$(R_s^2+\omega_e^2L_s^2)\matr I$.  The voltage boundary is a translated
circle.  This is the linear \gls{spmsm} limit; any plotted orientation is
meaningless because both eigenvalues coincide.

\paragraph{High speed.}
The inductive terms dominate, semi-axes scale approximately as
$1/\abs{\omega_e}$, and
\begin{equation}
 \ic\longrightarrow[-\psif/\Ld,\ 0]\trans.
\end{equation}
The feasible voltage region contracts about the $d$-axis current that cancels
PM flux.

\paragraph{Low speed.}
For $R_s>0$, $\Hpsm\to R_s^2\matr I$, $\ic\to0$, and the voltage boundary
approaches a resistance-limited circle of radius $U_{\max}/R_s$.  If both
$R_s=0$ and $\omega_e=0$, the stationary equations predict zero voltage for
every current and no finite voltage conic exists; current and thermal limits
then dominate.

\paragraph{IPMSM versus SPMSM.}
An \gls{ipmsm} has saliency, so unequal eigenvalues generally produce unequal
axes and reluctance torque bends the torque contours.  An ideal linear SPMSM
has equal inductances: its voltage boundary is circular and torque contours are
lines.  These conclusions concern the stated constant-inductance model; real
saturation can reintroduce anisotropy.

\subsection{Counterfactual checks}

\paragraph{What if the PM flux vanished?}
Then $\vect b=\vect0$, the voltage ellipse is centered at the origin, and only
reluctance torque remains.  The shape matrix is unchanged: excitation moves
the conic but does not create its anisotropy.

\paragraph{What if one eigenvalue became zero?}
The corresponding semi-axis formula would diverge.  The boundary would cease
to be a bounded ellipse and would become strip-like or degenerate depending on
the lower-order terms.  In the present model this requires rank loss of
$\Apsm$, reached at the joint limit $R_s=\omega_e=0$ for positive
inductances.

\paragraph{What if one eigenvalue were negative?}
The level set would have opposite curvature in its two principal directions
and would be hyperbolic.  This cannot happen for
$\Hpsm=\Apsm\trans\Apsm$, because a squared norm is never negative.  The
counterexample is useful: positive semidefiniteness is a structural consequence
of the voltage norm, not a numerical coincidence.

\paragraph{What if saturation destroys the quadratic model?}
Current-dependent inductances make the voltage map nonlinear, so one global
matrix can no longer classify the entire boundary.  Local Jacobians still
provide directional gains and local SVDs, while numerical level sets reveal
the global deformation.  The ellipse then becomes a reference geometry rather
than an exact law.

\begin{figure}[tb]
\centering
\begin{tikzpicture}
 \begin{axis}[axis main,width=.82\linewidth,height=6.5cm,
 xlabel={$i_d$ (normalized)},ylabel={$i_q$ (normalized)},
 xmin=-3.1,xmax=.8,ymin=-1.6,ymax=1.6,axis equal image]
  \addplot[voltage limit,BrickRed!55,domain=0:360,samples=160]
    ({-.45+2.15*cos(x)},{1.25*sin(x)});
  \addplot[voltage limit,BrickRed!78,domain=0:360,samples=160]
    ({-1.05+1.45*cos(x)},{.82*sin(x)});
  \addplot[voltage limit,BrickRed,domain=0:360,samples=160]
    ({-1.42+.95*cos(x)},{.52*sin(x)});
  \legend{low speed,medium speed,high speed}
 \end{axis}
\end{tikzpicture}
\caption{Increasing speed contracts the voltage ellipse and moves its center
toward the flux-cancellation current.}
\label{fig:parameter-change}
\end{figure}


For a representative set $R_s=\SI{0.08}{\ohm}$,
$L_d=\SI{0.8}{\milli\henry}$, $L_q=\SI{1.2}{\milli\henry}$,
$\psi_f=\SI{0.075}{\weber}$, and $U_{\max}=\SI{180}{\volt}$, the equations
above are sufficient to compute center, orientation, and axes without fitting
an ellipse to sampled points.  Numerical plots should be used to evaluate the
model, not to infer its mathematical class.

\subsection{Numerical verification of the derived geometry}

For lookup-table, measured, or \gls{fe} flux maps, analytic rearrangement may
be unavailable.  A robust verification samples the voltage residual on a mesh,
extracts the zero contour triangle by triangle, and compares it with the
analytic center, eigenvectors, and semi-axes in the linear case.  The checks
should include the residual $\abs{\norm{\voltage}-\Umax}$ and gradient
parallelism at claimed tangencies.  Delaunay contour extraction is useful for
irregular samples, but it does not replace the rank and definiteness proof: a
finite plotting window can hide an open or degenerate level set.
